% !TeX root = ../../Book5.tex
\chapter*{序言}
\addcontentsline{toc}{chapter}{序言}
\markboth{序言}{序言}

把一个正方形转过四分之一周，四个顶点换了位置，平面上的每个向量也跟着转动。前者给出顶点集的置换，后者给出一个线性变换；我们观察的是同一个群，却得到两种不同的描述。若只关心某个顶点去了哪里，置换已经足够。若想知道有没有保持不变的直线，或者几个作用能否拆成更小的部分，矩阵和向量空间更便于动手。这就是表示论的朴素想法：让一个代数对象作用在线性空间上，再由这些作用认识它。

\ProseParagraphBreak
这里的“线性”并不是把原来的问题简单改写成矩阵。两个不同的表示可能属于同一个群，一个表示也可能看不见群中的某些元素；换基后的矩阵虽然不同，描述的表示却可以相同。因此要研究的不只是矩阵本身，还有保持作用的线性映射、不变子空间和不可约分解。表示足够小的时候，人们可以列出矩阵逐项计算；表示多起来以后，便需要能够把许多计算合并的工具。

\ProseParagraphBreak
这门学问的一段来历很能说明它为何重视计算。1896年，Dedekind 在给 Frobenius 的信中提出群行列式及其分解问题；Frobenius 沿着这个问题发展出广义群特征标。当时的出发点并不是今天课本里“先给矩阵，再取迹”的讲法\cite[Section 4.11]{EtingofEtAl2011RepresentationTheory}。现代教材改从表示及其迹讲起，有了更清楚的结构，也没有丢掉那个最初的疑问：群的乘法如何在一个可以分解、可以计算的表达式里显出来？

\ProseParagraphBreak
第一编从有限群表示和群代数入手。先看正则表示、置换表示等熟悉的例子，再证明适当特征条件下的完全可约性。随后用特征标的正交关系计算不可约表示，用诱导和限制在群与子群之间传递信息，最后讨论双陪集、Clifford 理论和诱导定理。特征标表可以回答一些关于表示的问题，却不能代替群的全部结构；本编会区分表中能够读出的信息与仍需另证的事实。

\ProseParagraphBreak
第二编先到对称群。分拆、Young 图与 Specht 模，使“寻找不可约表示”变成一个兼有代数与组合意味的问题：一幅方格图既记录维数，也约束限制、诱导和张量运算。模特征的表示作为选读内容，提醒读者不要把复表示里的完全可约性任意搬到别的域上。之后转向箭图：几个点、几支箭头，看起来比群简单，表示却已经可以有丰富的分解情形。路径代数把箭头上的线性映射组织成模，维数向量和反射方法进一步揭示分类与图形之间的联系。

\ProseParagraphBreak
第三编讨论李代数。Sophus Lie 研究连续变换群和微分方程时，逐渐发展出研究无穷小变换的办法；这段历史也提醒我们，李代数虽然能够用矩阵括号举例，其问题最初并非为矩阵分类而设\cite[Section 2.10]{EtingofEtAl2011RepresentationTheory}。本编先建立定义、可解性、半单性和 \(\mathfrak{sl}_2\) 的表示，再以泛包络代数和 PBW 定理连接李代数与结合代数。Cartan 子代数、根系、Weyl 群及 Dynkin 图给出半单情形的组织方法；最高权理论和 Weyl 特征标公式则回答不可约表示怎样命名、维数与权重怎样计算。

\ProseParagraphBreak
第四编的三个选读章进一步讨论范畴 \(\mathcal O\)、李代数上同调、Hopf 代数、代数群和不变量理论。从上同调解释扩张，到用坐标环记录群运算，再到由不变多项式构造商，前几卷的模、正合性与交换代数在这里有了新的用途。将这些专题置于卷末，是为了使已有基础的读者能够继续深究，而不让初读有限群表示或最高权理论的人预先背负全部工具。

\ProseParagraphBreak
五篇附录各有侧重。组合代数与反射群延续 Young 表、计数和对称的讨论；量子群、簇代数与范畴化用小秩计算展示新的代数结构。李群、算术表示与几何表示论的简介说明，加入拓扑或几何条件以后，表示问题会怎样改变。最后两篇附录从微分算子和微分模出发，把线性方程的解空间与代数群联系起来。它们可以按兴趣选读；具体模型尽量展开计算，需要额外理论的一般结论则注明所用条件和出处。

\ProseParagraphBreak
不同分支各有自己的起点。读有限群表示，应先熟悉第一卷的群作用、共轭与商群，并掌握第二卷前七章关于线性空间、模和张量积的基本知识，完全可约性则在本卷建立。想读箭图，可以从本卷第九章及其两顶点、三顶点例子开始，再补前面关于模分解的工具；想读李代数，可先有线性算子、矩阵交换子和一些域的知识，不必等到交换代数或导出范畴全都读完。进入第四编的范畴 \(\mathcal O\) 和上同调时，再按章补第四卷的相关内容。

\ProseParagraphBreak
阅读时还要随时看清底域和特征。有限群的复表示可以借平均算子分解，模特征表示却可能留下不可分解的扩张；复半单李代数的最高权理论，也不能无条件照搬到正特征。“特征不整除群阶”或“特征为零”等假设决定着结论能否成立，须与所用结果一同记住。读者也可以试着寻找去掉条件后的第一个反例。

\ProseParagraphBreak
学习表示论不妨从小对象反复练习。给循环群和正方形的对称群写出作用矩阵，算出迹并比较共轭类；给一支箭头 \(V\to W\) 选不同秩的线性映射，看哪些表示同构；给 \(\mathfrak{sl}_2\) 的小维数模写出三个标准算子的作用。第一次计算是为得到答案，第二次则可以追问答案中哪些部分与基的选择无关。等这一步想清楚了，再读一般定理，符号通常会亲近许多。

\ProseParagraphBreak
从有限群的特征标表走到根系与最高权，符号和工具已经丰富了许多。最值得保留的仍是开头正方形所提示的两种眼光：既看见一个对称在具体对象上做了什么，也看见许多不同作用背后共享的结构。希望读者终能带着这两种眼光，在有限群、箭图与李代数之间找到自己的问题；也愿本卷的计算与证明，能帮助读者把最初的疑问写成一个可以认真研究的数学问题。

\bigskip
\begin{flushright}
R. EverStarine

2026年9月26日
\end{flushright}
